\section{Setup and conventions}
\label{sec:setup}

\paragraph{Model.}
Spatially flat FRW background, general relativity with $\Lambda=0$, signature $(-+++)$, and a
single perfect radiation fluid
\begin{equation}
T^{\mu\nu}=(\rho+P)u^\mu u^\nu+P g^{\mu\nu},\qquad P=\rho/3,\qquad \rho=\alpha\,\Tphys^4,
\label{eq:fluid}
\end{equation}
with $\alpha$ constant and $u^\mu u_\mu=-1$.  The second relation \emph{defines} the physical
temperature $\Tphys$; no chemical potential or particle-number assumption is made.  Test photons
are emitted isotropically with a Planck spectrum of temperature $T_*$ on the hypersurface
$\Tphys=T_*$, then propagate collisionlessly without backreaction.

\paragraph{Background and units.}
Conformal time $\eta$, $\dd s^2=a^2(\eta)\,\hat g_{\mu\nu}\dd x^\mu\dd x^\nu$.  In radiation
domination $a\propto\eta$; we fix the normalisation $a(\eta)=\eta$, so that $\mathcal H\equiv a'/a=1/\eta$
and the Friedmann equation $3\mathcal H^2=8\pi G a^2\bar\rho$ reads
\begin{equation}
8\pi G\,a^4\bar\rho=3,\qquad\text{i.e.}\qquad 8\pi G\,\bar\rho=3/\eta^4 .
\label{eq:friedmann}
\end{equation}
Background temperature $\bar T(\eta)=\bar T_i\,\eta_i/\eta$.  The \emph{conformal} fluid variables
\begin{equation}
\hat\rho\equiv a^4\rho,\qquad \hat P\equiv a^4P,\qquad \hat T\equiv a\,\Tphys,\qquad
\hat u^\mu\equiv a\,u^\mu\quad(\hat g_{\mu\nu}\hat u^\mu\hat u^\nu=-1),
\label{eq:hatvars}
\end{equation}
are constant in the background ($\hat{\bar\rho}=3/8\pi G$, $\hat{\bar T}=$ const) and are the natural
variables of Sec.~\ref{sec:observable}.

\paragraph{Fourier conventions and statistics.}
\begin{equation}
f(\vx)=\int\!\frac{\dd^3k}{(2\pi)^3}\,e^{i\vk\cdot\vx}f_{\vk},\qquad
\langle\zeta_{\vk}\zeta_{\vk'}\rangle=(2\pi)^3\delta^3(\vk+\vk')\,P(k),\qquad
\langle\zeta^2(\vx)\rangle=\int\!\frac{\dd^3k}{(2\pi)^3}P(k)=\frac1{2\pi^2}\int k^2\dd k\,P(k).
\label{eq:fourier}
\end{equation}
Products in position space become convolutions, $ (fg)_{\vk}=\int\dd^3q/(2\pi)^3\,f_{\vq}\,g_{\vk-\vq}$,
and $\partial_i\to ik_i$, $\lap\to-k^2$, $\ilap\to-1/k^2$.  $\langle\cdot\rangle$ denotes the ensemble
average over the Gaussian field $\zeta_1$ (Sec.~\ref{sec:primordial}); $\sky{\cdot}$ denotes the
solid-angle average over the observer's sky.

\paragraph{Perturbative expansion.}
Every perturbation $X\in\{\Psi,\Phi,\omega_i,h_{ij},\delta,v^i,\ldots\}$ is expanded as
\begin{equation}
X=\eps X_1+\tfrac{\eps^2}{2}X_2+\tfrac{\eps^3}{6}X_3+O(\eps^4),
\label{eq:expansion}
\end{equation}
and the same convention is used for every derived quantity ($S$, $\Lambda$, $m$, $\Theta$, $\zeta_{\rm nl}$).
For an equation $E[X]=0$ we write $E^{(n)}\equiv n!\,[\eps^n]E$, so that $X_n$ enters $E^{(n)}$ linearly with
unit weight.  The bookkeeping parameter $\eps$ is the amplitude of $\zeta_1$; the final expansion
parameter is $\eps^2\langle\zeta_1^2\rangle$.  Products of lower-order quantities are never dropped.

\paragraph{Poisson gauge.}
The conformal metric is parametrised \emph{linearly} (no exponentials),
\begin{equation}
\hat g_{00}=-(1+2\Psi),\qquad \hat g_{0i}=\omega_i,\qquad \hat g_{ij}=(1-2\Phi)\delta_{ij}+h_{ij},
\qquad \partial_i\omega_i=0,\quad h_{ii}=0=\partial_ih_{ij},
\label{eq:poisson}
\end{equation}
with $\omega_i$ and $h_{ij}$ starting at order $\eps^2$ (no primordial vector or tensor modes,
Sec.~\ref{sec:primordial}).  Indices on $\omega_i$, $h_{ij}$, $v^i$ are raised and lowered with $\delta_{ij}$.
The fluid is parametrised by
\begin{equation}
\rho=\bar\rho(\eta)\,(1+\delta),\qquad 1+\delta_T\equiv(1+\delta)^{1/4}=\hat T/\hat{\bar T},\qquad
\hat u^\mu=\hat u^0\,(1,v^i),
\label{eq:fluidvars}
\end{equation}
with $\hat u^0$ fixed by the normalisation $\hat g_{\mu\nu}\hat u^\mu\hat u^\nu=-1$ as a series in $\eps$.
Irrotationality of $v^i$ is \emph{not} imposed (Sec.~\ref{sec:second}).  The Einstein equations are
$E_{\mu\nu}\equiv G_{\mu\nu}[g]-8\pi G\,T_{\mu\nu}=0$ with $g=\eta^2\hat g$ and
$8\pi G\,T_{\mu\nu}=(3/\eta^4)(1+\delta)[\tfrac43u_\mu u_\nu+\tfrac13g_{\mu\nu}]$; the conservation laws
$\nabla_\mu T^{\mu\nu}=0$ are implied and are used as checks.
Dictionary to the exponential parametrisation of \cite{BMR1,BMR2,MZ}
($\hat g_{00}=-e^{2\Phi_{\rm e}}$, $\hat g_{ij}=e^{-2\Psi_{\rm e}}\delta_{ij}+\ldots$, note the swapped names):
$\Phi_{\rm e}^{(2)}=\Psi_2-2\Psi_1^2$, $\Psi_{\rm e}^{(2)}=\Phi_2+2\Phi_1^2$ at second order.

\paragraph{Directions, rays and times.}
$\nhat$ is the direction in which the observer looks (unit vector with respect to the observer's
orthonormal triad, Sec.~\ref{sec:observer}); the photon propagation direction at the observer is
$\ehat=-\nhat$.  The background ray through the observer event $(\bar\eta_o,\bar\vx_o)$ is
\begin{equation}
\bar\vx(\eta)=\bar\vx_o+\nhat\,\chi(\eta),\qquad \chi(\eta)\equiv\bar\eta_o-\eta,\qquad
\frac{\dd\bar\vx}{\dd\eta}=-\nhat,
\label{eq:bgray}
\end{equation}
and the total derivative along it acting on a field is $D\equiv\partial_\eta-\nhat\cdot\nabla$.  Four times are
kept distinct: the \emph{initial} time $\eta_i$ at which initial data are posed
($\bar T(\eta_i)=T_i$); the background \emph{emission} time $\eta_*$ defined by $\bar T(\eta_*)=T_*$, and the
actual emission time $\eta_e(\nhat)=\eta_*+\delta\eta_e(\nhat)$ along each ray; the \emph{observation} time
$\eta_o=\bar\eta_o+\delta\eta_o$ of the observer (direction independent, fixed by proper time); and the
dummy time $\eta$ along the ray.  The ordering is $0<\eta_i<\eta_*<\bar\eta_o$.
The angular quantities are
\begin{equation}
a^{(r)}_{\ell m}=\int\dd\Omega_{\nhat}\,Y^*_{\ell m}(\nhat)\,\Theta^{(r)}(\nhat),\qquad
C^{(rs)}_\ell=\frac1{2\ell+1}\sum_{m=-\ell}^{\ell}\big\langle a^{(r)}_{\ell m}a^{(s)*}_{\ell m}\big\rangle_c ,
\label{eq:cl_def}
\end{equation}
with $Y_{\ell m}$ carrying the Condon--Shortley phase, and the target is
\begin{equation}
C_\ell=\eps^2C^{(11)}_\ell+\eps^4\Big[\tfrac14C^{(22)}_\ell+\tfrac13\Re C^{(13)}_\ell\Big]+O(\eps^6),\qquad \ell\ge2 .
\label{eq:cl_target}
\end{equation}

\paragraph{Relation to existing work.}
No complete $O(\eps^4)$ (one-loop) computation of the temperature power spectrum exists in the literature:
the only $(22)$-type temperature calculation \cite{Tomita} is for dust in synchronous gauge (ISW terms only) and
explicitly leaves out the $(13)$ term for lack of third-order metric perturbations; third-order temperature
anisotropies have been studied on large scales in matter domination \cite{DAmico,Gao}; second-order Boltzmann
codes \cite{SONG,FKP} and the $(22)$-type $B$-mode spectrum \cite{FidlerBmode} exist, as do the second-order
Sachs--Wolfe analyses \cite{Boubekeur,HuangVernizzi,BMR1,BMR2,PUB} which, ``following standard practice'',
drop the observer terms; only \cite{MZ,CPV} keep them, for long-wavelength modes.  The lensed $C_\ell$ at
$O(C^\phi)$ \cite{Hu2000,LewisChallinor} is the $(22)+(13)$ structure restricted to lensing and serves as our
control case.  The radiation-era second-order vector and tensor sectors are treated in \cite{Lu} and
\cite{KT,Baumann,Espinosa}; general second-order perturbation theory is reviewed in \cite{MalikWands}.
